\documentclass[12pt]{article}
\usepackage[margin=1in]{geometry}

% Start of preamble
%==========================================================================================%
% Required to support mathematical unicode
\usepackage[warnunknown, fasterrors, mathletters]{ucs}
\usepackage[utf8x]{inputenc}

\usepackage[dvipsnames,table,xcdraw]{xcolor}
\usepackage{hyperref} 
\hypersetup{
colorlinks=true,
linkcolor=blue,
filecolor=magenta,
urlcolor=cyan,
pdfpagemode=FullScreen
}

% Standard mathematical typesetting packages
\usepackage{amsmath,amssymb,amscd,amsthm,amsxtra, pxfonts}
\usepackage{mathtools,mathrsfs,dsfont,xparse}

% Symbol and utility packages
\usepackage{cancel, textcomp}
\usepackage[mathscr]{euscript}
\usepackage[nointegrals]{wasysym}
\usepackage{apacite}

% Extras
\usepackage{physics}  
\usepackage{tikz-cd} 
\usepackage{microtype}
\usepackage{enumitem}
\usepackage{titling}
\usepackage{graphicx}

% Fancy theorems due to @intuitively on discord
\usepackage{mdframed}
\newmdtheoremenv[
backgroundcolor=NavyBlue!30,
linewidth=2pt,
linecolor=NavyBlue,
topline=false,
bottomline=false,
rightline=false,
innertopmargin=10pt,
innerbottommargin=10pt,
innerrightmargin=10pt,
innerleftmargin=10pt,
skipabove=\baselineskip,
skipbelow=\baselineskip
]{mytheorem}{Theorem}

\newenvironment{theorem}{\begin{mytheorem}}{\end{mytheorem}}

\newtheorem{corollary}{Corollary}
\newtheorem{lemma}{Lemma}

\newtheoremstyle{definitionstyle}
{\topsep}%
{\topsep}%
{}%
{}%
{\bfseries}%
{.}%
{.5em}%
{}%
\theoremstyle{definitionstyle}
\newmdtheoremenv[
backgroundcolor=Violet!30,
linewidth=2pt,
linecolor=Violet,
topline=false,
bottomline=false,
rightline=false,
innertopmargin=10pt,
innerbottommargin=10pt,
innerrightmargin=10pt,
innerleftmargin=10pt,
skipabove=\baselineskip,
skipbelow=\baselineskip,
]{mydef}{Definition}
\newenvironment{definition}{\begin{mydef}}{\end{mydef}}

\newtheorem*{remark}{Remark}

\newtheorem*{example}{Example}

% Common shortcuts
\def\mbb#1{\mathbb{#1}}
\def\mfk#1{\mathfrak{#1}}

\def\bN{\mbb{N}}
\def \C{\mbb{C}}
\def \R{\mbb{R}}
\def\bQ{\mbb{Q}}
\def\bZ{\mbb{Z}}
\def \cph{\varphi}
\renewcommand{\th}{\theta}
\def \ve{\varepsilon}
\newcommand{\mg}[1]{\| #1 \|}

% Often helpful macros
\newcommand{\floor}[1]{\left\lfloor#1\right\rfloor}
\newcommand{\ceil}[1]{\left\lceil#1\right\rceil}
\renewcommand{\qed}{\hfill\qedsymbol}
\renewcommand{\P}{\mathbb P\qty}
\newcommand{\E}{\mathbb{E}\qty}
\newcommand{\Cov}{\mathrm{Cov}\qty}
\newcommand{\Var}{\mathrm{Var}\qty}

% Sets
\usepackage{braket}

\graphicspath{{/}}
\usepackage{float}

\newcommand{\SET}[1]{\Set{\mskip-\medmuskip #1 \mskip-\medmuskip}}

% End of preamble
%==========================================================================================%

% Start of commands specific to this file
%==========================================================================================%

%==========================================================================================%
% End of commands specific to this file

\title{CSE 446 HW1}
\date{\today}
\author{Rohan Mukherjee}

\begin{document}
    \maketitle
    \begin{enumerate}[leftmargin=\labelsep]
        \item \begin{enumerate}
            \item Given the $x_i$, the probability of seeing the $x_i$ in that order is given by the following likelihood function:
            \begin{align*}
                \mathscr{L}(\lambda) = \prod_{i=1}^n e^{-\lambda} \frac{\lambda^{x_i}}{x_i!}
            \end{align*}
            Taking the log gives,
            \begin{align*}
                \log \mathscr{L}(\lambda) = \sum_{i=1}^n -\lambda + x_i \log \lambda - \log x_i!
            \end{align*}
            The derivative of this function is just
            \begin{align*}
                \dv{\lambda} \log \mathscr{L}(\lambda) = \sum_{i=1}^n -1 + \frac{x_i}{\lambda}
            \end{align*}
            Setting this to 0 yields $\lambda = \frac{1}{n} \sum_{i=1}^n x_i$.

            \item From this, we can see that the numerical estimate for $\lambda$ for the first 5 games is just
            \begin{align*}
                \lambda = \frac15 (3+7+5+0+2) = 3.4
            \end{align*}
            Thus, the probability that the Reign score exactly 4 goals in the next game is just
            \begin{align*}
                \P(4 \mid \lambda) = e^{-3.4}\frac{3.4^4}{4!} \approx 0.168
            \end{align*}

            \item Given that Reign actually scored 8 goals in their last game, the updated $\lambda$ is just
            \begin{align*}
                \lambda = \frac16 (3+7+5+0+2+8) = 4.1667
            \end{align*}
            Now, the probability that Reign scores 5 goals in their next game is just
            \begin{align*}
                \P(5 \mid \lambda) = e^{-4.1667}\frac{4.1667^5}{5!} \approx 0.156
            \end{align*}
        \end{enumerate}
    \end{enumerate}
\end{document}